% GENERATED FILE. Edit canonical lessons, then run npm run generate:books.
% canonical-source-hash: fnv1a64:ef9e2a07adc901f1
% canonical-lessons: GU-C284-adhar, GU-C284-adhar-rakhe-che, GU-C284-hot-to-avat, GU-W54-six-lines, GU-R284-depends-and-would

\chapter{It depends, and what would have been}
\label{ch:gu-it-depends-and-what-would-have-been}
\hlchaptermodality{284}

\begin{chapteropening}
\textbf{By the end of this chapter:} \emph{I can say what something depends on (… par ādhār rākhe che), say what would happen or would have happened (jo … hot, to … āvat), and write six connected lines about a trip.}

Complete GU-R284-depends-and-would: say what something depends on, and what would have been.
\end{chapteropening}

\section[\texorpdfstring{\gu{આધાર} (ādhār)}{આધાર (ādhār)}]{\gu{આધાર} (ādhār) — a support; what something rests on}

\label{lesson:GU-C284-adhar}



\begin{quote}
Before the new one: say \textquotedblleft{}or else\textquotedblright{}. (\emph{nahī̃tar}.) Then \textquotedblleft{}to keep\textquotedblright{}. (\textbf{\gu{રાખવું}}.)
\end{quote}

\subsection*{You'll want to know: \gu{આધાર}}
\textbf{\gu{આધાર}} — \emph{ādhār} — \textquotedblleft{}a support; what something rests on\textquotedblright{}. It is masculine: \textbf{\gu{મારો} \gu{આધાર}}, \textquotedblleft{}my support\textquotedblright{}.

\begin{quote}
\textbf{\gu{મારાં} \gu{માતા}-\gu{પિતા} \gu{મારો} \gu{આધાર} \gu{છે}.}
\end{quote}

\begin{quote}
\emph{mārā̃ mātā-pitā māro ādhār che.} — \textquotedblleft{}My parents are my support.\textquotedblright{}
\end{quote}

A person can be a \textbf{\gu{આધાર}}, and so can a fact that a plan rests on. The next lesson builds \textquotedblleft{}it depends\textquotedblright{} on that second sense.

\begin{cousinweb}
\textbf{\gu{આધાર}} (\emph{ādhār}) is a tatsama: Sanskrit \emph{ādhāra}, \textquotedblleft{}support, prop\textquotedblright{}, and \textquotedblleft{}the power of holding up\textquotedblright{} (Monier-Williams), from \emph{ā-} and the root \emph{dhṛ-}, \textquotedblleft{}to hold\textquotedblright{}.
\end{cousinweb}

\subsection*{Your turn}
Say these aloud:

\begin{itemize}
\raggedright
  \item \emph{ādhār}, then \emph{māro ādhār}
  \item \textquotedblleft{}My parents are my support\textquotedblright{} — \emph{mārā̃ mātā-pitā māro ādhār che}
\end{itemize}

\begin{tcolorbox}[breakable,colback=teal!4,colframe=teal!35!black,title={Before you move on}]
What does \textbf{\gu{આધાર}} mean? (\textquotedblleft{}A support; what something rests on\textquotedblright{}.) Its gender? (Masculine: \textbf{\gu{મારો} \gu{આધાર}}.)
\end{tcolorbox}
\section[\texorpdfstring{… \gu{પર} \gu{આધાર} \gu{રાખે} \gu{છે} (… par ādhār rākhe che)}{… પર આધાર રાખે છે (… par ādhār rākhe che)}]{… \gu{પર} \gu{આધાર} \gu{રાખે} \gu{છે} (… par ādhār rākhe che) — it depends on …}

\label{lesson:GU-C284-adhar-rakhe-che}



\begin{quote}
Say \textquotedblleft{}my support\textquotedblright{}. (\emph{māro ādhār}.) Then \textquotedblleft{}if it rains\textquotedblright{}. (\emph{jo varsād paḍśe}.)
\end{quote}

\begin{grammarlens}[title={Grammar Lens: it keeps its support on …}]
\begin{quote}
\textbf{\gu{તે} \gu{હવામાન} \gu{પર} \gu{આધાર} \gu{રાખે} \gu{છે}.}
\end{quote}

\begin{quote}
\emph{te havāmān par ādhār rākhe che.} — \textquotedblleft{}It depends on the weather.\textquotedblright{}
\end{quote}

Word for word: \textquotedblleft{}it keeps its support on the weather\textquotedblright{}. \textbf{\gu{પર}} (\emph{par}), \textquotedblleft{}on\textquotedblright{}, comes after its word, as \textbf{\gu{સુધી}} does. The thing that depends is the subject of \textbf{\gu{રાખે} \gu{છે}}: \textbf{\gu{બધું} \gu{સમય} \gu{પર} \gu{આધાર} \gu{રાખે} \gu{છે}} (\emph{badhũ samay par ādhār rākhe che}), \textquotedblleft{}everything depends on the time\textquotedblright{}. A person takes \textbf{\gu{મારા}}, \textbf{\gu{તમારા}}: \textbf{\gu{તે} \gu{તમારા} \gu{પર} \gu{આધાર} \gu{રાખે} \gu{છે}} (\emph{te tamārā par ādhār rākhe che}), \textquotedblleft{}it depends on you\textquotedblright{}.
\end{grammarlens}

\subsection*{You'll want to know: \textquotedblleft{}it depends\textquotedblright{}, and then}
\begin{quote}
\textbf{\gu{અમે} \gu{જઈશું} \gu{કે} \gu{નહીં}? \gu{તે} \gu{વરસાદ} \gu{પર} \gu{આધાર} \gu{રાખે} \gu{છે}.}
\end{quote}

\begin{quote}
\emph{ame jaīśũ ke nahī̃? te varsād par ādhār rākhe che.} — \textquotedblleft{}Will we go or not? It depends on the rain.\textquotedblright{}
\end{quote}

\textquotedblleft{}It depends\textquotedblright{} is seldom the end of an answer. Say on what, then follow with \textbf{\gu{જો}}: \textbf{\gu{જો} \gu{વરસાદ} \gu{નહીં} \gu{પડે}, \gu{તો} \gu{અમે} \gu{જઈશું}} (\emph{jo varsād nahī̃ paḍe, to ame jaīśũ}).

\subsection*{Your turn}
Say these aloud:

\begin{itemize}
\raggedright
  \item \textquotedblleft{}Everything depends on the time\textquotedblright{} — \emph{badhũ samay par ādhār rākhe che}
  \item \textquotedblleft{}It depends on you\textquotedblright{} — \emph{te tamārā par ādhār rākhe che}
\end{itemize}

\begin{tcolorbox}[breakable,colback=teal!4,colframe=teal!35!black,title={Before you move on}]
\textquotedblleft{}It depends on the weather\textquotedblright{}? (\textbf{\gu{તે} \gu{હવામાન} \gu{પર} \gu{આધાર} \gu{રાખે} \gu{છે}}.) Where does \textbf{\gu{પર}} go? (After its word.) \textquotedblleft{}On you\textquotedblright{}? (\textbf{\gu{તમારા} \gu{પર}}.)
\end{tcolorbox}
\section[\texorpdfstring{\gu{જો} … \gu{હોત}, \gu{તો} … \gu{આવત} (jo … hot, to … āvat)}{જો … હોત, તો … આવત (jo … hot, to … āvat)}]{\gu{જો} … \gu{હોત}, \gu{તો} … \gu{આવત} (jo … hot, to … āvat) — what would have been}

\label{lesson:GU-C284-hot-to-avat}



\begin{quote}
Say \textquotedblleft{}it depends on the rain: if it doesn't rain, we'll go\textquotedblright{}. (\emph{te varsād par ādhār rākhe che: jo varsād nahī̃ paḍe, to ame jaīśũ}.) Then \textquotedblleft{}I have time\textquotedblright{}, the A2 way. (\textbf{\gu{મારી} \gu{પાસે} \gu{સમય} \gu{છે}}.)
\end{quote}

\begin{grammarlens}[title={Grammar Lens: the -\gu{ત} form}]
\begin{quote}
\textbf{\gu{જો} \gu{મારી} \gu{પાસે} \gu{સમય} \gu{હોત}, \gu{તો} \gu{હું} \gu{આવત}.}
\end{quote}

\begin{quote}
\emph{jo mārī pāse samay hot, to hũ āvat.} — \textquotedblleft{}If I had had time, I would have come.\textquotedblright{}
\end{quote}

The stem with \textbf{-\gu{ત}} (\emph{-t}) is \textquotedblleft{}would\textquotedblright{}: the condition did not hold, so the result did not happen. Both halves take it: \textbf{\gu{હોત}}, from \textbf{\gu{હોવું}}, and \textbf{\gu{આવત}}, from \textbf{\gu{આવવું}}. It can also be about now, a thing that is not so: \textquotedblleft{}if I had time, I would come\textquotedblright{}.

\textquotedblleft{}If not\textquotedblright{} puts \textbf{\gu{ન}} before it: \textbf{\gu{જો} \gu{ટ્રેન} \gu{મોડી} \gu{ન} \gu{હોત}, \gu{તો} \gu{અમે} \gu{વહેલા} \gu{પહોંચત}} (\emph{jo ṭren moḍī na hot, to ame vahelā pahõcat}), \textquotedblleft{}if the train hadn't been late, we would have arrived early\textquotedblright{}.
\end{grammarlens}

\subsection*{You'll want to know: one form for everyone}
\textbf{-\gu{ત}} stays the same for every person, man or woman: \textbf{\gu{હું} \gu{જાત}}, \textbf{\gu{તે} \gu{જાત}}, \textbf{\gu{અમે} \gu{જાત}}, with the same \emph{jāt} each time. It is no past, so a verb with an object keeps a plain doer, with no \textbf{-\gu{એ}}:

\noindent
\begin{tabularx}{\linewidth}{@{}>{\raggedright\arraybackslash}X>{\raggedright\arraybackslash}X@{}}
\toprule
\textbf{plain past} & \textbf{would have} \\
\midrule
\textbf{\gu{મેં} \gu{કામ} \gu{કર્યું}} & \textbf{\gu{હું} \gu{કામ} \gu{કરત}} \\
\textbf{\gu{તેણે} \gu{કહ્યું}} & \textbf{\gu{તે} \gu{કહેત}} \\
\bottomrule
\end{tabularx}

\subsection*{Your turn}
Say these aloud:

\begin{itemize}
\raggedright
  \item \textquotedblleft{}If the train hadn't been late, we would have arrived early\textquotedblright{} — \emph{jo ṭren moḍī na hot, to ame vahelā pahõcat}
  \item \textquotedblleft{}If I had time, I would do the work\textquotedblright{} — \emph{jo mārī pāse samay hot, to hũ kām karat}
\end{itemize}

\begin{tcolorbox}[breakable,colback=teal!4,colframe=teal!35!black,title={Before you move on}]
\textquotedblleft{}I would have come\textquotedblright{}? (\textbf{\gu{હું} \gu{આવત}}.) Does it change for a woman? (No.) \textquotedblleft{}I would do the work\textquotedblright{}: \textbf{\gu{મેં}} or \textbf{\gu{હું}}? (\textbf{\gu{હું} \gu{કામ} \gu{કરત}}.)
\end{tcolorbox}
\section[(six lines)]{Six lines about a trip}

\label{lesson:GU-W54-six-lines}



\begin{quote}
The A2 composition was a practical note. This task asks for six lines that tell a story: a trip by train, bus or boat, real or invented.

Pick one reader: Arun, \textbf{\gu{અરુણ}}, or Mira, \textbf{\gu{મીરા}}. Mira's Sunday at the village is your model of order.
\end{quote}

\subsection*{What to know first}
A told story is not a list. It has a \textbf{scene}, with \textbf{-\gu{તો} \gu{હતો}} or \textbf{\gu{જ્યારે} … \gu{ત્યારે}}, and \textbf{events} in the plain past, kept in order with \textbf{\gu{પહેલાં}}, \textbf{\gu{ત્યાર} \gu{પછી}} and \textbf{\gu{છેવટે}} and joined with \textbf{-\gu{ઈને}}. It has a \textbf{reason}, with \textbf{\gu{કારણ} \gu{કે}}, \textbf{-\gu{ને} \gu{લીધે}} or \textbf{\gu{એટલે}}, and a \textbf{feeling}, with \textbf{\gu{મને} \gu{બીક} \gu{લાગી}} or \textbf{\gu{હું} \gu{ખુશ} \gu{હતો}}.

A good one ends on what might have been: \textbf{\gu{જો} …, \gu{તો} … -\gu{ત}}.

\subsection*{Writing — connected composition}
On paper, in Gujarati script, write \textbf{six lines} to your reader about the trip, one line for each of these:

\begin{enumerate}
\raggedright
  \item the scene: when it was, and what was going on as it began, with \textbf{-\gu{તો} \gu{હતો}} or \textbf{\gu{જ્યારે} … \gu{ત્યારે}};
  \item the first thing that happened, with \textbf{\gu{પહેલાં}}, and a \textbf{-\gu{ઈને}} form;
  \item what happened next, with \textbf{\gu{ત્યાર} \gu{પછી}}, and one thing that went wrong: \textbf{\gu{ખોવાઈ} \gu{ગઈ}} or \textbf{\gu{છૂટી} \gu{ગઈ}};
  \item why, with \textbf{\gu{કારણ} \gu{કે}}, \textbf{-\gu{ને} \gu{લીધે}} or \textbf{\gu{એટલે}};
  \item how it ended and how you felt, with \textbf{\gu{છેવટે}} and \textbf{\gu{મને} \gu{દુઃખ} \gu{થયું}} or \textbf{\gu{હું} \gu{ખુશ} \gu{હતો}};
  \item what might have been, with \textbf{\gu{જો} … \gu{હોત}, \gu{તો} … -\gu{ત}}, and one comparison, with \textbf{\gu{કરતાં}} or \textbf{\gu{સૌથી}}.
\end{enumerate}

There is no model in this lesson. Use only words you already write. The work is choosing the right past for each line, and joining them.

\begin{tcolorbox}[breakable,colback=teal!4,colframe=teal!35!black,title={Before you move on}]
Now read your lines as you read Mira's Sunday. Mark every verb \textbf{s} for scene, \textbf{e} for event. Line one should be mostly \textbf{s}; lines two and three mostly \textbf{e}.

Then three checks. A verb with an object in the plain past needs \textbf{-\gu{એ}} on a name: \textbf{\gu{અરુણે}}, \textbf{\gu{મીરાએ}}. \textbf{\gu{દુઃખ}} has \textbf{\gu{ઃ}} after \textbf{\gu{દુ}}, and \textbf{\gu{સૌથી}} has \textbf{\gu{ૌ}} on its \textbf{\gu{સ}}. \textbf{\gu{ઝડપી}} opens on \textbf{\gu{ઝ}}.

Last, read line six alone. Do both halves take \textbf{-\gu{ત}}: \textbf{\gu{હોત}}, then \textbf{\gu{આવત}} or \textbf{\gu{જાત}}?
\end{tcolorbox}
\section[(review)]{Review — it depends, and what would have been}

\label{lesson:GU-R284-depends-and-would}



\begin{quote}
Say \textquotedblleft{}if it rains, we'll stay at home\textquotedblright{}. (\emph{jo varsād paḍśe, to ame ghare rahīśũ}.) Then \textquotedblleft{}or else\textquotedblright{}. (\emph{nahī̃tar}.)
\end{quote}

\subsection*{Your turn}
Say these aloud:

\begin{itemize}
\raggedright
  \item My parents are my support. What do you think?
  \item It depends on the weather. Everything depends on the time.
  \item Will we go or not? It depends on the rain. If it doesn't rain, we'll go.
  \item If you're tired, stay at home. If I miss the train, I'll come by bus: the bus is slower than the train, but it's cheaper.
  \item If I had had time, I would have come.
  \item If the train hadn't been late, we would have arrived early.
\end{itemize}

\begin{tcolorbox}[breakable,colback=teal!4,colframe=teal!35!black,title={Before you move on}]
\textquotedblleft{}It depends on you\textquotedblright{}? (\textbf{\gu{તે} \gu{તમારા} \gu{પર} \gu{આધાર} \gu{રાખે} \gu{છે}}.) \textquotedblleft{}I would have come\textquotedblright{}, said by a woman? (\textbf{\gu{હું} \gu{આવત}}: the same.) \textquotedblleft{}I would do the work\textquotedblright{}: with \textbf{-\gu{એ}} or without? (Without: \textbf{\gu{હું} \gu{કામ} \gu{કરત}}.)
\end{tcolorbox}
